\chapter{Comparar sem Fingir: Humano, Ferramenta, Agente e Equipa}
\label{ch:comparar}

«A equipa resolveu mais» pode significar mais texto, mais tentativas ou mais tarefas realmente aceites. Antes de comparar, precisamos de decidir a unidade e conservar o que cada condição recebeu. Caso contrário, o vencedor pode ter sido escolhido quando escrevemos o enunciado.

\section{Quatro condições, quatro perguntas}
Uma pessoa sem IA fornece uma referência de trabalho humano naquele contexto. Uma pessoa com IA mede uma interação, incluindo aprendizagem, revisão e integração. Um agente único mede um sistema com o seu conjunto de ferramentas e orçamento. Uma equipa mede outro sistema, com custos de coordenação e oportunidades diferentes.

Não substituas a primeira condição por uma persona «humano» produzida por um LLM. Isso seria uma simulação pelo modelo, não um participante humano. Também não uses a taxa de testes do código para preencher a linha de qualidade de respostas do modelo.

\begin{table}[htbp]
\small
\begin{tabularx}{\textwidth}{@{}lXX@{}}
\toprule
Condição & O que registar & Fonte de aceitação \\
\midrule
Humano & Tempo total, experiência e ferramentas permitidas & Resultado verificado por critérios \\
Humano com IA & Pedidos, revisão, correções e tempo & Mesmos critérios, identidade oculta quando viável \\
Agente único & Versão, contexto, chamadas e falhas & Contratos e revisão de conteúdo \\
Equipa & Todos os anteriores e coordenação & Resultado final e erros dos componentes \\
\bottomrule
\end{tabularx}
\caption{Condições propostas; esta edição não apresenta medições com participantes humanos.}
\end{table}

\section{Controlar a comparação que nos interessa}
Se a pergunta é qualidade com orçamento igual, fixa o orçamento e deixa cada arquitetura usá-lo dentro das regras. Se a pergunta é custo para atingir uma qualidade, define o limiar e mede recursos. Se a pergunta é experiência de uso, tarefas e participantes precisam de representar esse objetivo.

Não é obrigatório tornar tudo igual; é obrigatório explicitar o que difere. Uma equipa com três chamadas e uma chamada única não têm o mesmo orçamento de inferência. Essa comparação pode ser útil, desde que os resultados mostrem os custos ao lado do acerto.

Usa tarefas separadas de desenvolvimento e uma rubrica escrita antes. Baralha ordem quando fizer sentido, conserva saídas e falhas, e procura efeitos por tipo de caso. Se uma mesma pessoa executa duas condições, prática e memória podem influenciar a segunda. Um piloto pequeno serve para melhorar o protocolo; raramente sustenta uma afirmação populacional forte.

\section{O que foi implementado nesta edição}
O controlador V3 tem modo único e equipa de dois especialistas com síntese. O adaptador por omissão é determinístico. Há testes de orçamento concorrente, contratos, snapshots e falhas parciais. O adaptador HTTP aceita apenas servidor local, usa timeout e recusa redirecionamentos.

O piloto planeado usa quatro pedidos sintéticos, duas seeds e um modelo Qwen pequeno. A condição de fluxo usa regras; as condições único e equipa usam chamadas distintas. As respostas esperadas e os prompts foram guardados antes da tentativa de inferência. O objetivo é categoria e fonte, sem resposta enviada a cliente nem operação de escrita.

As 32 chamadas planeadas com LLM \textbf{não foram executadas}: o runtime recusou o modelo local incompleto e não foi possível obter outra cópia através da ligação disponível. O registo em \texttt{experiencias/equipas/README.md} conserva erro, tamanho e hashes. Não existe uma tabela de ganhos de equipa a preencher com números inventados.

Os registos reais da segunda edição são outra experiência: 52 inferências, com configuração e tarefas próprias, no capítulo~\ref{ch:avaliacao}. Podem ser reanalisados, mas não completam retroativamente este piloto. Manter essa separação é parte do resultado metodológico.

\section{Um contraexemplo offline útil}
Executar o mesmo raciocínio determinístico em vários papéis não cria informação nova. O ensaio offline associado usa os quatro pedidos e aplica o mesmo stub nos modos de fluxo, único e equipa. Os outputs e as contagens podem ser comparados; as chamadas reais de modelo são zero.

Se os três modos falham a mesma paráfrase, a síntese não teve acesso a uma capacidade diferente que resolvesse a falha. Isso é um contraexemplo de software construído, não um estudo que prove inutilidade de todas as equipas. A sua utilidade é tornar visível a diferença entre multiplicar chamadas e acrescentar um mecanismo de correção.

O ensaio offline foi executado em 04-11-2026. A verificação repete o registo guardado em \nolinkurl{experiencias/equipas/resultados}. Nesta simulação, a linha «fluxo» usa o mesmo ramo determinístico da chamada única. A linha «equipa» acrescenta duas etapas no controlador. Assim, esta comparação observa chamadas de componentes e acerto do stub nas condições dadas.

\begin{lstlisting}
python3 experiencias/equipas/offline.py --verificar
\end{lstlisting}

\begin{table}[htbp]
\small
\centering
\begin{tabular}{@{}lccc@{}}
\toprule
Modo & Casos corretos & Chamadas de componentes & Chamadas por acerto \\
\midrule
Fluxo & 2 em 4 & 4 & 2 \\
Único & 2 em 4 & 4 & 2 \\
Equipa & 2 em 4 & 12 & 6 \\
\bottomrule
\end{tabular}
\caption{Resultados offline. Cada modo fez zero chamadas a modelos.}
\label{tab:comparacao-offline}
\end{table}

Confere também os mesmos casos lado a lado. Os três modos falharam E1 e E2 e acertaram E3 e E4. A diferença emparelhada de acerto entre equipa e único é zero nestes quatro pedidos. A razão de chamadas por acerto é $12/2=6$ na equipa e $4/2=2$ no único. Se não houver acertos, deixa a razão indefinida e apresenta as contagens. A conclusão fica circunscrita a estes pedidos e a este adaptador.

\section{Custos, erros graves e distribuição dos resultados}
Mostra tarefas aceites, encaminhamentos, erros críticos, chamadas, tokens quando existentes e tempo observado. Para tarefas abertas, uma nota de clareza não deve compensar uma violação de autorização. Podes usar critérios de exclusão e depois comparar qualidade entre as saídas admissíveis.

Apresenta denominadores e dependências. Com duas sementes nos mesmos quatro pedidos, agrupa as repetições por pedido. Se selecionares a melhor tentativa, explica a seleção e contabiliza todas as chamadas. Para custo, regista também as tarefas interrompidas. Para tempo, mede a duração até ao resultado aceite, incluindo revisão e fila se fizerem parte da utilização.

\begin{promptlabbox}{Um relatório que deixa linhas honestamente vazias}
\textbf{Objetivo e contexto:} preparar a comparação completa e executar a parte offline. Requisitos: Python 3.10+ para o controlador; participantes e modelo íntegro apenas para extensões futuras.

\textbf{Procedimento:} escreve tarefa, rubrica, dados e orçamento; executa os testes e o ensaio offline documentado; preenche apenas resultados observados. Para humano e humano com IA, deixa «não executado» até haver um estudo consentido e efetivo. Para LLM, usa pasta nova e integridade confirmada.

\textbf{Resultado esperado e validação:} o relatório permite distinguir medição, exemplo construído e protocolo. Outra pessoa consegue repetir a parte executável sem API paga.

\textbf{Falhas possíveis:} simular participantes e chamá-los humanos, misturar ensaios, omitir revisão ou anunciar superioridade global com base num teste de quatro casos.
\end{promptlabbox}

Se voltares a correr o piloto com um modelo íntegro, guarda o novo registo ao lado do ensaio offline. Mantém em separado as condições com pessoas, os resultados do modelo e as verificações do controlador. Cada uma responde a uma pergunta diferente.
